% 由粗到细的金字塔光流教学构造：真实位移8.6像素，三层逐层减半；方块边长按每层减半绘制。宽约108mm。
\begin{tikzpicture}[x=1mm,y=1mm,font=\figurefont\small,
  lvl/.style={draw=FigureBlue,fill=FigureBlueFill,line width=.6pt},
  link/.style={-{Stealth[length=2mm]},draw=FigureStroke,line width=.7pt},
  lab/.style={font=\figurefont\footnotesize,text=FigureMuted,align=center}]
  \draw[lvl] (0,9) rectangle (8,15);
  \draw[lvl] (26,6) rectangle (42,18);
  \draw[lvl] (60,0) rectangle (92,24);
  \node at (4,19) {第2层};
  \node at (34,22) {第1层};
  \node at (76,28) {第0层（原图）};
  \draw[-{Stealth[length=1.4mm]},FigureRed,line width=.9pt] (3,12) -- (5.2,12);
  \draw[-{Stealth[length=1.4mm]},FigureRed,line width=.9pt] (31,12) -- (35.3,12);
  \draw[-{Stealth[length=1.4mm]},FigureRed,line width=.9pt] (71,12) -- (79.6,12);
  \draw[link] (9,12) -- (25,12) node[midway,above,lab] {位移$\times2$};
  \draw[link] (43,12) -- (59,12) node[midway,above,lab] {位移$\times2$};
  \node[lab] at (4,3) {初值$0$\\估计$2.1$};
  \node[lab] at (34,0) {初值$4.2$\\细化为$4.3$};
  \node[lab] at (76,-6) {初值$8.6$\\细化为$8.6$};
\end{tikzpicture}
